% Regression: a selector may name a derived carrier.  Its resolved boundary
% exit, rather than an invented grid cell, determines routed crossings.
\documentclass{standalone}
\usepackage{tenkz}
\begin{document}
\tenkzkernel
\begin{tenkz}[rows={wire}, cols=3, bonds=none]
  \tn[at=(1,1), name=A]{A}
  \tn[at=(1,3), name=B]{B}
  \tn[at=midway A and B, name=M, ports={90:virtual}]{M}
  \tnwire[
    kind=string,
    name=north-route,
    route={n of M},
    crossing=over
  ]{open}{open}
\end{tenkz}
\begin{tenkz}[rows={wire}, cols=3, bonds=none]
  \tn[at=(1,1), name=C]{C}
  \tn[at=(1,3), name=D]{D}
  \tn[at=midway C and D, name=N, ports={0:virtual}]{N}
  \tnwire[
    kind=string,
    name=off-face-route,
    route={n of N},
    crossing=over
  ]{open}{open}
\end{tenkz}
\begin{tenkz}[rows={wire}, cols=3, bonds=none]
  \tn[at=(1,1), name=E]{E}
  \tn[at=(1,3), name=F]{F}
  \tnwire[name=carrier]{E.0}{F.180}
  \tn[at=on carrier 0.5, name=O, ports={90:virtual}]{O}
  \tnwire[
    kind=string,
    name=onwire-route,
    route={n of O},
    crossing=over
  ]{open}{open}
\end{tenkz}
\end{document}
